% ======================================================================
% 08-conclusion.tex
% ======================================================================
\section{Conclusion}
\label{sec:conclusion}

The project did not become reusable because it acquired a package layout or a
test count. It became reusable because every failure left a guard, every guard
left an artifact, and every artifact was written where the next session would
look. The five mechanisms are the visible residue of that process: one reference
implementation with a parity-tested fast path, manifests that make a run
self-describing, gates that turn physical assumptions into go/no-go decisions, a
statistical harness that makes a single seed embarrassing, and an artifact-first
runner that treats the run rather than the process as the unit.

The scientific results that survive are modest and specific: a 1D baseline
reported as a median with a 7$\times$ spread; an FNO repaired from worse than
zeros to a 3\% median; two 2D policies 1.2\% apart on different well sets with a
400-scenario certificate; and a 3D paper preset reproduced on a third machine.
The results that do not survive are just as useful: $\tau$-ramping, a three-well
policy, a reported CVaR tail, and a planned GPU speedup. Keeping both lists is
the practice.

The AI contribution was real but narrower than the marketing: it made
verification, repetition and record-keeping cheap enough to be habitual. The
human contribution was to decide what counted as evidence. The repository
contribution was to remember. A project that keeps all three---and writes down
which is which---can be reused by someone who was not there.

\paragraph{Code and data availability.}
The repository contains the solvers, configs, tests, run manifests and artifacts
referenced in this article. The reproducibility entry points are
\code{python tests/run\_tests.py}, \code{python SCRIPTS/build\_tex.py
--check-only}, and the per-section commands in Appendix~\ref{app:repro}. The
article source, evidence ledger and source map are in \code{ARTICLE/}.
